\subsection{Il problema del trasferimento dello stato}

Siano $\mathcal T_h^n$ e $\mathcal T_h^{n+1}$ due triangolazioni con la stessa connettività ma coordinate nodali differenti, e siano $V_h^n$ e $V_h^{n+1}$ i corrispondenti spazi $P_1$. In generale, $V_h^n\neq V_h^{n+1}$: una rilocazione istantanea della mesh richiede quindi di trasferire lo stato $u_h^n$ nel nuovo spazio discreto. Due operatori sono stati implementati e confrontati come diagnostica, in accordo con la letteratura sul trasferimento conservativo tra mesh \citep{zzhang2006moving,alauzet2010p1}.

\subsection{Trasferimento puntuale \texorpdfstring{$P_1$}{P1}}

Il remapping puntuale valuta la funzione lineare definita sulla vecchia mesh nei nodi della nuova mesh. Se il nuovo nodo $\bm x_i^{n+1}$ appartiene all'elemento $K_e^n$ con coordinate baricentriche $\lambda_{e,a}$, il valore trasferito è
\begin{equation}
U_i^{n+1}=
u_h^n(\bm x_i^{n+1})
=\sum_{a=1}^{3}\lambda_{e,a}(\bm x_i^{n+1})U_{e,a}^n.
\label{eq:p1_remap}
\end{equation}
L'operatore è locale e relativamente economico, ma non coincide in generale con un'identità tra gli spazi discreti. La sua applicazione ripetuta può introdurre errori di rappresentazione e modificare l'integrale del campo.

\subsection{Proiezione variazionale \texorpdfstring{$L^2$}{L2}}

Una seconda possibilità definisce $u_h^{n+1}\in V_h^{n+1}$ imponendo
\begin{equation}
\int_\Omega v_h^{n+1}u_h^{n+1}\,d\bm x
=
\int_\Omega v_h^{n+1}u_h^n\,d\bm x
\qquad \forall v_h^{n+1}\in V_h^{n+1}.
\label{eq:l2_projection}
\end{equation}
In forma matriciale:
\begin{equation}
M^{n+1}\bm U^{n+1}=\bm b^{n\rightarrow n+1},
\end{equation}
con $M_{ij}^{n+1}=(\phi_j^{n+1},\phi_i^{n+1})_\Omega$ e
$b_i^{n\rightarrow n+1}=(u_h^n,\phi_i^{n+1})_\Omega$. Se la base di prova soddisfa la partizione dell'unità e gli integrali sono valutati esattamente sullo stesso dominio, scegliendo come funzione test la costante unitaria si ottiene l'identità
\begin{equation}
\int_\Omega u_h^{n+1}\,d\bm x=\int_\Omega u_h^n\,d\bm x.
\label{eq:l2_integral}
\end{equation}
Questa proprietà riguarda il trasferimento $L^2$ idealizzato; in un calcolo numerico la sua accuratezza dipende da quadratura, localizzazione geometrica e soluzione del sistema.

\subsection{Formulazione Arbitrary Lagrangian--Eulerian}

Nella campagna principale il movimento non viene applicato come salto istantaneo seguito da un remapping. Si introduce una mappa di coordinate
\begin{equation}
\bm x=\bm\chi(\widehat{\bm x},t),\qquad
\bm w(\widehat{\bm x},t)=
\left.\frac{\partial\bm\chi}{\partial t}\right|_{\widehat{\bm x}},
\label{eq:ale_map}
\end{equation}
dove $\widehat{\bm x}$ è la coordinata di riferimento e $\bm w$ la velocità della mesh. Per una funzione sufficientemente regolare, la regola della catena dà
\begin{equation}
\left.\partial_tu\right|_{\widehat{\bm x}}
=
\left.\partial_tu\right|_{\bm x}+\bm w\cdot\nabla u.
\label{eq:ale_chain}
\end{equation}
La PDE \eqref{eq:pde} può pertanto essere espressa, nella descrizione mobile, come
\begin{equation}
\left.\partial_tu\right|_{\widehat{\bm x}}
+(\bm b-\bm w)\cdot\nabla u
-\nu\Delta u=f.
\label{eq:ale_strong}
\end{equation}
La formulazione discreta opera con una matrice di massa dipendente dal tempo e con i contributi convettivi relativi alla velocità fisica e alla velocità della griglia. La variazione della geometria è trattata all'interno del passo ALE; non viene applicato un operatore separato di remapping dello stato nella pipeline Stage D.

Per esplicitare l'operatore effettivamente assemblato, si indichino con $M^n$ e $M^{n+1}$ le matrici di massa sulle configurazioni ai due estremi del passo, con $C_{\mathrm{phys}}^{n+1}$ il contributo convettivo fisico, con $K^{n+1}$ quello diffusivo e con $F^{n+1}$ il vettore di forzante. Il contributo cinematico della mesh è assemblato alla configurazione intermedia e definito da
\begin{equation}
C_{\mathrm{mesh}}^{n+1/2}
=-C_{\mathrm{adv}}^{w,n+1/2}-M_{\mathrm{div}w}^{n+1/2},
\label{eq:ale_mesh_operator}
\end{equation}
dove, schematicamente, $(C_{\mathrm{adv}}^{w})_{ij}=\int_{\Omega^{n+1/2}}\varphi_i\,\bm w\cdot\nabla\varphi_j\,d\bm x$ e $(M_{\mathrm{div}w})_{ij}=\int_{\Omega^{n+1/2}}\varphi_i\varphi_j\nabla\cdot\bm w\,d\bm x$. Il passo backward Euler implementato è quindi rappresentabile come
\begin{equation}
\begin{aligned}
\left[\Delta t^{-1}M^{n+1}
+C_{\mathrm{mesh}}^{n+1/2}
+C_{\mathrm{phys}}^{n+1}+K^{n+1}\right]\bm U^{n+1}
&=\Delta t^{-1}M^n\bm U^n\\
&\quad+\bm F^{n+1}.
\end{aligned}
\label{eq:ale_discrete_step}
\end{equation}
con imposizione forte delle condizioni di Dirichlet sui gradi di libertà di bordo. La scrittura distingue l'operatore cinematico, valutato a metà passo, dai termini fisici assemblati sulla configurazione finale; non implica che tutti i contributi siano valutati a un unico tempo intermedio.

\paragraph{Interpolazione temporale del movimento.}
A ogni finestra di adattamento, di ampiezza $H_a=10\Delta t=0.05$, la GNN propone uno spostamento e il filtro restituisce lo spostamento ammissibile $\Delta\bm X_{\mathrm{safe}}$. Tale spostamento viene distribuito linearmente sui dieci passi della finestra. La velocità di mesh è quindi costante a tratti:
\begin{equation}
\bm w=\frac{\Delta\bm X_{\mathrm{safe}}}{H_a}.
\label{eq:mesh_velocity}
\end{equation}
La topologia resta invariata; il solutore ALE nativo aggiorna la geometria e risolve il sistema FEM nel dominio mobile.

\subsection{Legge geometrica di conservazione discreta}

La \emph{Discrete Geometric Conservation Law} (\DGC) esprime la compatibilità tra variazione della geometria della mesh e discretizzazione temporale. In una formulazione ALE, la matrice di massa e gli operatori dipendono dalla configurazione; il termine associato alla velocità di mesh deve essere coerente con tale dipendenza. Una verifica fondamentale consiste nella preservazione del campo costante quando l'evoluzione è ridotta al solo trasporto geometrico, senza dinamica fisica che ne modifichi il valore \citep{farhat2001the,guillard2000on,formaggia2004stability}.

La verifica viene effettuata separando l'identità geometrica element-wise, il residuo algebrico assemblato, l'errore di free-stream e il residuo del sistema lineare. Nella verifica indipendente Stage E1 su una mesh strutturata $16\times16$ sono stati misurati un residuo geometrico nativo element-wise massimo $8.69\times10^{-17}$, un residuo algebrico in norma infinito $4.63\times10^{-16}$ e un errore free-stream in norma infinito $4.10\times10^{-13}$. Il residuo lineare è stato pari a $1.13\times10^{-12}$, contro una tolleranza GMRES effettiva registrata pari a $2.27\times10^{-11}$. Questi risultati riguardano il test indipendente E1 e non devono essere confusi con l'audit più ampio Stage D, nel cui log storico un test free-stream non superava una soglia più stretta di $10^{-12}$. I residui geometrici documentano la compatibilità della configurazione verificata; non implicano, da soli, la conservazione della massa fisica per la PDE forzata. Nel problema corrente le condizioni di Dirichlet sono omogenee e stazionarie, come in \eqref{eq:homogeneous_dirichlet}; la massa varia comunque per effetto della forzante e del flusso diffusivo normale al bordo, secondo \eqref{eq:continuous_mass_balance}. La D-GCL, la discrepanza degli incrementi di massa e il residuo integrale del bilancio PDE sono pertanto tre oggetti distinti.

\subsection{Implementazione numerica}

Il nucleo numerico è implementato in C++20 con deal.II, utilizzando le strutture di triangolazione, gli spazi agli elementi finiti, gli oggetti di integrazione e le matrici/solutori lineari della libreria \citep{arndt2021dealii}. I binding pybind11 espongono le funzionalità native a Python. La GNN, la gestione dei dataset e l'orchestrazione restano in Python. Le formule geometriche del filtro e alcuni assemblaggi ALE/remapping sono implementati direttamente nel core C++; non si attribuisce quindi a deal.II l'implementazione di ogni singolo kernel. L'architettura evita duplicazioni Python per le operazioni di mesh e FEM.
